%! Polynomial Functions and Rates of Change — Unit 2, Lesson 2.2 — Homework (Answer Key)
\documentclass[10pt]{article}
\usepackage{precalculus-article}
\usepackage{precalculus-key}   % pulls in -boxes; do NOT also load -boxes
\newcommand{\TallMath}[1]{$\displaystyle #1\rule[-1.4em]{0pt}{3.2em}$}

\begin{document}

\pageheader{Unit 2, Lesson 2.2}{Homework --- Answer Key}

\vspace{0.12in}

\begin{enumerate}[itemsep=10pt]

\item The graph of a polynomial $f$ on the restricted domain $[0, 5]$ is
      shown.

\begin{center}
\begin{tikzpicture}[x=0.8cm, y=0.5cm]
  \draw[linegray, very thin, step=1] (0,0) grid (5,5);
  \draw[->, thick] (0,0) -- (5.6,0) node[below right] {\small $x$};
  \draw[->, thick] (0,0) -- (0,5.6) node[above] {\small $f(x)$};
  \foreach \x in {1,2,3,4,5} \node[below] at (\x,0) {\small \x};
  \foreach \y in {1,2,3,4,5} \node[left] at (0,\y) {\small \y};
  \draw[very thick, plum] plot [smooth] coordinates
    {(0,1) (0.5,3.2) (1,5) (1.5,3.9) (2,2.2) (2.5,0.7) (3,0)
     (3.5,0.5) (4,1.7) (4.5,3.0) (5,4)};
  \fill[goldacc] (1,5) circle (2.7pt);
  \fill[goldacc] (3,0) circle (2.7pt);
\end{tikzpicture}
\end{center}

\begin{enumerate}[label=(\alph*), itemsep=4pt]
  \item Increasing on \ans{$(0,1)$} and \ans{$(3,5)$} ; decreasing on
        \ans{$(1,3)$} .
  \item Local maximum value \ans{5} at $x = $ \ans{1} ;
        local minimum value \ans{0} at $x = $ \ans{3} .
  \item Both endpoints are included in the domain. The endpoint $x = 0$ is a
        local \ans{minimum} and the endpoint $x = 5$ is a local
        \ans{maximum} .
  \item Global maximum value \ans{5} ; global minimum value
        \ans{0} .
\end{enumerate}

\boxguard[18]
\item The table gives values of a polynomial $q$.

\smallskip
\renewcommand{\arraystretch}{1.5}
\begin{tabular}{|c|c|c|c|c|c|c|}
\hline
$x$ & 0 & 1 & 2 & 3 & 4 & 5 \\
\hline
$q(x)$ & 2 & 9 & 14 & 17 & 22 & 31 \\
\hline
\end{tabular}
\smallskip

AROCs over $[0,1], [1,2], [2,3], [3,4], [4,5]$: \quad
\ans{7} , \ans{5} , \ans{3} , \ans{5} , \ans{9}

\begin{enumerate}[label=(\alph*), itemsep=4pt]
  \item Does $q$ have a local maximum or minimum on $[0,5]$? \quad
        \textbf{Yes} / \ans{No} \quad Why? \ans{every AROC is positive}
  \item The AROCs decrease and then increase, so a point of inflection lies
        near $x = $ \ans{3} .
  \item There the graph changes from concave \ans{down} to concave
        \ans{up} .
\end{enumerate}

\item A polynomial $p$ has exactly four distinct real zeros:
      $x = -3$, $x = 0$, $x = 2$, and $x = 5$.

\begin{enumerate}[label=(\alph*), itemsep=4pt]
  \item At least how many local extrema must $p$ have? \ans{3}
  \item Explain your answer in one sentence.

        \vspace{0.05in}
        \ansline{There are three gaps between consecutive zeros, and each gap forces at least one turn.}
\end{enumerate}

\item Decide whether each statement \textbf{must} be true. Circle one and give
      a one-line reason.

\begin{enumerate}[label=(\alph*), itemsep=6pt]
  \item A polynomial of degree 6 has either a global maximum or a global
        minimum. \quad \ans{True} / \textbf{False}

        Reason: \ans{Even degree --- both ends run the same way, so outputs are bounded on one side.}

  \item A polynomial of degree 5 has a global maximum. \quad
        \textbf{True} / \ans{False}

        Reason: \ans{Odd degree --- one end rises without bound, so no output is the greatest.}

  \item A polynomial with three distinct real zeros has at least two local
        extrema. \quad \ans{True} / \textbf{False}

        Reason: \ans{Two gaps between consecutive zeros, each forcing at least one turn.}
\end{enumerate}

\item The average rates of change of a polynomial $p$ over the five
      consecutive intervals $[0,1], [1,2], [2,3], [3,4], [4,5]$ are
      $8,\ 5,\ 3,\ 4,\ 6$. Which statement must be true?

\begin{enumerate}[label=\Alph*., itemsep=3pt]
  \item $p$ has a local minimum somewhere on $(0, 5)$.
  \item The graph of $p$ has a point of inflection on $(0, 5)$. \quad \textcolor{keyred}{\textbf{$\leftarrow$ correct: AROCs fall then rise}}
  \item $p$ has a global maximum at $x = 5$.
  \item $p$ is decreasing on part of $(0, 5)$.
\end{enumerate}

\end{enumerate}

\vspace{0.1in}

\boxguard
\begin{extensionbox}
A polynomial $r$ has degree 5 and exactly five distinct real zeros. Use both
of today's counting rules --- ``at least one turn between consecutive zeros''
and ``at most $n-1$ turns for degree $n$'' --- to pin down exactly how many
local extrema $r$ has.

\vspace{0.05in}
\ansline{Five zeros make four gaps, so $r$ has at least 4 local extrema. Degree 5 allows at most}
\ansline{$5 - 1 = 4$. The two bounds meet, so $r$ has exactly 4 local extrema.}
\end{extensionbox}

\vspace{0.1in}

\boxguard
\begin{notesbox}{Preview of Lesson 2.3}
Today the zeros of a polynomial were only landmarks --- the fence posts that
forced a turn between them. In Lesson 2.3, \textit{Polynomial Functions and
Real Zeros}, the zeros become the main event: how to find them from a factored
form, and what the graph does when it reaches one.
\end{notesbox}

\end{document}
